\section{Аналіз помилок застосування асиметричної криптографії}

\subsection{Ситуація 1. Неправильне використання padding у RSA}

\textbf{Умова:} RSA використовується без рекомендованої схеми padding або із застарілою схемою, непридатною для відповідного сценарію.

\textbf{Помилка:}

\textbf{Можливі наслідки:}

\textbf{Спосіб усунення:}

\subsection{Ситуація 2. Повторне використання nonce в ECDSA}

\textbf{Умова:} Під час створення двох різних ECDSA-підписів було повторно використано однаковий nonce.

\textbf{Помилка:}

\textbf{Можливі наслідки:}

\textbf{Спосіб усунення:}

\subsection{Ситуація 3. Передбачуваний генератор випадкових чисел}

\textbf{Умова:} Криптографічні ключі генеруються з використанням передбачуваного генератора випадкових чисел.

\textbf{Помилка:}

\textbf{Можливі наслідки:}

\textbf{Спосіб усунення:}

\subsection{Ситуація 4. Відсутність перевірки відкритого ключа}

\textbf{Умова:} Система отримує відкритий ключ від іншої сторони та використовує його без перевірки його належності очікуваному учаснику взаємодії.

\textbf{Помилка:}

\textbf{Можливі наслідки:}

\textbf{Спосіб усунення:}

\subsection{Ситуація 5. Повторне використання ключової пари}

\textbf{Умова:} Одна й та сама пара ключів використовується одночасно для різних криптографічних призначень без аналізу сумісності таких способів використання.

\textbf{Помилка:}

\textbf{Можливі наслідки:}

\textbf{Спосіб усунення:}

\subsection{Загальний висновок щодо помилок застосування}